\paragraph{One loop about the conjugate-coordinate family at fixed \(\gamma>0\): the fermion form of the finite-\(N\) members (conditional; bounded no-go).}
Keep Definition~(A) with \(h_0\) of signature \((9,1)\) in place of \(\eta\), the mass term \(S_\gamma\) with \(\gamma>0\), the fermion content of~\cite{IKKT1997}, supplied as there, and the family \(\bar T^I=L^I{}_aT^a\), \(T^a=M^{ac_0}\), \(H=L^{\mathrm T}h_0L=c\,\eta\), of the preceding paragraph. The one-loop term about a member is the Gaussian integral over the bosonic and ghost fluctuations in the background Feynman gauge times the Pfaffian of the fermion form \(\operatorname{Tr}\psi^{\mathrm T}C\Gamma_I[\bar T^I,\psi]\) at the member, as in the one-loop formula of~\cite{Steinacker2023}, whose \(\mathrm i\epsilon\) is the regulator of Definition~(A). A member whose four-plane \(W\) is spacelike is carried, by an isometry of \(h_0\) in its identity component, which acts invertibly on the spinor index, to \(\bar T^a=\sqrt c\,M^{ac_0}\) along four spacelike basis directions of a basis in which \(h_0=\eta\).

\emph{Operators of the compact members.} Let the five-index algebra be \(\mathfrak{so}(5)\), with \(W\) spacelike and \(c>0\), so that \(\gamma=3c\), and use that form. Write \(C_5\) and \(C_4\) for the Casimirs of \(\mathfrak{so}(5)\) and of its four-index subalgebra \(\mathfrak{so}(4)\) in the adjoint action on the matrices, normalised to \(4\) and \(3\) on the vector, and \(C_4^{\mathrm{tot}}\) for the Casimir of \(\mathfrak{so}(4)\) acting on the matrices and on a vector or spinor index together. The bosonic Hessian in this gauge is then \(c\,(C_5-C_4^{\mathrm{tot}})\) on the four directions of \(W\) and \(\pm c\,(C_5-C_4-3)\) on each of the six others, with the sign of \(h_0\) there; the ghost operator is \(c\,(C_5-C_4)\); and the square of \(\Gamma^a[\bar T_a,\cdot\,]\) is \(c\,(C_5-C_4^{\mathrm{tot}}+\tfrac32)\) on each Weyl spinor of \(\mathfrak{so}(4)\), of which the chiral spinor of \(\mathrm{SO}(9,1)\) contains four of each chirality (exact, with spinor generators obeying the relations of the \(M^{ab}\); checked on explicit matrices for the representations \(n=1,2\) below). The ten bosons, sixteen fermion components and two ghosts weight the scale \(c\) by \(5-4-1=0\), so that, apart from zero modes, the modulus of the one-loop term does not depend on \(\gamma\).

\emph{Fermion zero modes.} Label the irreducible components of the matrices under \(\mathfrak{so}(5)\) by \((m_1,m_2)\), \(m_1\ge m_2\ge0\), both integers or both half-integers, in the orthogonal basis, with \(C_5=m_1(m_1+3)+m_2(m_2+1)\), and those of \(\mathfrak{so}(4)\) by \((k_1,k_2)\), with \(j_{\mathrm L}=\tfrac12(k_1+k_2)\), \(j_{\mathrm R}=\tfrac12(k_1-k_2)\) and \(C_4=2j_{\mathrm L}(j_{\mathrm L}+1)+2j_{\mathrm R}(j_{\mathrm R}+1)\). The Gelfand--Tsetlin branching \(m_1\ge k_1\ge m_2\ge|k_2|\) contains \((k_1,k_2)=(m_1,m_2)\), and this component coupled to a spinor of type \((\tfrac12,0)\) contains total spin \(j_{\mathrm L}+\tfrac12\), on which \(C_4^{\mathrm{tot}}=2(j_{\mathrm L}+\tfrac12)(j_{\mathrm L}+\tfrac32)+2j_{\mathrm R}(j_{\mathrm R}+1)=C_5+\tfrac32\); likewise \((m_1,-m_2)\) with \(j_{\mathrm R}+\tfrac12\) for \((0,\tfrac12)\). On these subspaces the Hermitian operator \(\Gamma^a[\bar T_a,\cdot\,]\) vanishes. In a non-trivial finite-dimensional representation the generators are linearly independent and span a copy of the adjoint \((1,1)\) among the traceless matrices. So at every finite \(N\) the fermion form of every non-zero member has a kernel, of dimension at least \(32\) per copy of the adjoint, and \(\operatorname{Pf}\mathcal M(\bar T)=0\). By the preceding paragraph these are all the non-zero members at finite \(N\). The fermion form also has a kernel at \(N=\infty\) for every non-trivial unitary representation of \(\mathrm{Spin}(5)\), each being a direct sum of finite-dimensional ones, and on a product with a finite factor \(K\) in the transverse directions, where these modes tensored with the identity of \(K\) stay in the kernel. For \(V_n=\mathrm{Sym}^n\mathbb C^4\), which is irreducible under \(\mathfrak{so}(5)\cong\mathfrak{sp}(4)\) and has \(N=\binom{n+3}3\), the matrices are \(\bigoplus_{0\le m_2\le m_1\le n}(m_1,m_2)\), each with multiplicity one, these are the only zero modes, and the kernel has dimension \(\tfrac43n(n+4)(n^2+4n+7)\), that is \(80\) and \(304\) for \(n=1,2\), as on explicit matrices (Racah--Speiser and the branching for \(n\le20\)); random Hermitian configurations of four matrices give no kernel (sampled). The ratio of this dimension to the \(16(N^2-1)\) fermion components tends to zero like \(3/n^2\), and its ratio to \(N\), which is proportional to the symplectic volume, grows like \(8n\). Zero modes of the same Casimir type occur for fuzzy \(S^4_N\) with five matrices in the Euclidean model, where the sector containing them is left out of the one-loop potential~\cite{Steinacker2015S4} (its Eq.~(3.28) and Section~4).

\emph{Consequences for the compact members.} Since \(\operatorname{Pf}\mathcal M(\bar T)=0\) is a property of the integrand, the fermionic factor of the one-loop term vanishes at every finite-\(N\) member, whatever the regulator. Each Yukawa vertex absorbs at most two zero modes, so the fermionic integral about such a member starts at an order in the fluctuations of at least half the dimension of the kernel, and the contribution of the member is not of one-loop form. Its order in \(g_{\mathrm{YM}}/\gamma\) is not stated, since the bosonic Hessian has, besides the twenty-four zero modes of the rotations that move \(W\), further zero eigenvalues (\(6\), \(142\) and \(277\) for \(n=1,2,3\)) whose lifting is not computed. It is also indefinite, with \(34\), \(120\) and \(340\) negative eigenvalues on the traceless fluctuations along the four directions of \(W\) (the compact check of Section~\ref{sec:status-F} counts four more from the trace) and \(10\), \(19\) and \(19\) on each spacelike transverse direction for \(n=1,2,3\); in Definition~(A) negative eigenvalues change the phase of the one-loop term and not its modulus, so the sign of the fluctuation spectrum does not distinguish the moduli of the members' contributions.

\emph{Non-compact members and normalisation.} The operator identities above are identities in the enveloping algebra and hold for \(\mathfrak{so}(4,1)\supset\mathfrak{so}(3,1)\), with the Casimirs formed with \(\eta\). The non-compact members act, however, on infinite-dimensional unitary representations, whose components under the five-index algebra are not finite-dimensional, so the branching argument does not apply; whether their fermion form has a normalisable kernel is not decided here. By the zero modes, the members of this family with a non-vanishing one-loop term are at most those with a non-compact five-index algebra, which exist only for \(N\to\infty\). At fixed \(\gamma>0\) and \(h_0\) of signature \((9,1)\) the identity of the preceding paragraph admits, besides \(\mathfrak{so}(5)\), only the five-index algebra \(\mathfrak{so}(4,1)\) with \(W\) timelike, among them the conjugate-coordinate member, whose kinetic tensor is Lorentzian on \(m\neq0\); the base-coordinate member requires \(\gamma<0\) there. If the fermion forms of these members are non-degenerate, the supplied fermion content, and not the bosonic action, separates them from the compact members at one loop; the Lorentzian kinetic tensor is then fixed by the supplied sign of \(\gamma\), \(h_0\) and fermion content, not derived, and is not identified with the emergent metric. A comparison of the bosonic factors per unit volume, by contrast, depends on the regulator. The leading term in which the mass spoils the cancellation between the species is proportional to \(\gamma\operatorname{Tr}(K+\mathrm i\epsilon)^{-1}\), and its real part enters the modulus. For \(K=-k_0^2+|\vec k|^2\) on flat momenta, the real part of \(\int\mathrm d^4k\,(K+\mathrm i\epsilon)^{-1}\) is \(0\) after Wick rotation. Over the ellipsoid \(k_0^2/a^2+|\vec k|^2\le\Lambda^2\) it is \(2\pi^2\Lambda^2a/(1+a^2)\), against \(2\pi^2\Lambda^2a/(1+a)\) for \(K=k_0^2+|\vec k|^2\) over the same region (exact; numerical check), and Definition~(A) supplies no such cut-off at \(N=\infty\).

These statements are for Definition~(A) with the supplied mass, fermion content and family, one loop about the members and the background Feynman gauge. For the compact members they hold for finite-dimensional representations, exactly at every finite \(N\), and need no limit. For the non-compact members the order is \(N\to\infty\) first at fixed \(\gamma>0\), then \(\epsilon\to0\), as in the preceding paragraph. Their fermion form, the existence of that limit, the cut-off of their bosonic factor and the contributions of saddles outside the family are not treated. Nothing here derives the family, the mass or the fermion content, so Assumption~\ref{ass:lorentzian} is not discharged, and Problem~E remains open and PROGRAMME-CONDITIONAL.
